\section{Static Analysis}
\label{sec:analysis}

In this section, we present two static analysis techniques:
%
origin and capability analysis (\cref{sec:analysis:oricap}) and error source
analysis (\cref{sec:analysis:error}).
%
We also discuss the reasons for classifying certain streams as unsupported,
which prevents them from being replaced with Rust std streams in the subsequent
transformation phase (\cref{sec:analysis:unsupported}).
%
% While we analyze Rust code generated by C2Rust, the code examples in this
% section are written in C because C2Rust-generated code calls libc functions and
% retains C-style coding patterns.

\subsection{Origin and Capability Analysis}
\label{sec:analysis:oricap}

This analysis aims to determine the origins and capabilities of each
\emph{location} storing a stream, where a location is a variable or a struct
field of type \C{FILE*}.
%
The analysis results allow the transformation to assign a single Rust type to
each location.
%
Since each location's type is fixed throughout the entire program, the analysis
is flow-insensitive;
%
it computes the set of all origins possibly stored in each location and the set
of all capabilities required by the operations performed on each location.
%
We formally define a location $l$ as follows:
%
\[
  %
  \small
  %
  x\in\textrm{Variables}\qquad
  %
  t\in\textrm{Structs}\qquad
  %
  f\in\textrm{Fields}\qquad
  %
  l\ ::=\ x\ |\ t.f
  %
\]

In the analysis, an origin is either $\ostdin$, $\ostdout$, $\ostderr$,
$\ofile$, or $\opipe$, and a capability is either $\cread$, $\cbufread$,
$\cwrite$, $\cseek$, or $\cclose$.
%
Although closing a stream is not represented by a trait in Rust std, we treat it
as a capability because the transformation must decide whether to use an owning
or borrowing type based on whether the location closes the stream.

We now define the syntax of a program being analyzed, focusing only on the key
features relevant to the analysis:
%
\[
  %
  \small
  %
  \begin{array}{@{}r@{~}r@{~}l@{}}
    %
    p &::=& x\ |\ e.f \\
    %
    e &::=& p\ |\ {\tt stdin}\ |\ {\tt stdout}\ |\ {\tt stderr}\ |\ {\tt fopen}()\ |\
    {\tt popen}() \\
    %
    s &::=& p=e\ |\ {\tt fread}(p)\ |\ {\tt fscanf}(p)\ |\ {\tt fwrite}(p) \\
    %
    &|& {\tt fseek}(p)\ |\ {\tt fclose}(p)\ |\ {\tt pclose}(p) \\
    %
  \end{array}
  %
\]
%
A \emph{path} $p$ is either a variable or a field projection, representing a
left value of an assignment or an argument to a function call.
%
An \emph{expression} $e$ is either a path, standard input/output/error, a file
open, or a pipe construction, representing a right value of an assignment.
%
A \emph{statement} $s$ is either an assignment or a call to an API function that
takes a stream.
%
A program is a sequence of statements;
%
we ignore control structures as the analysis is flow-insensitive.

We finally describe the \emph{constraints} generated by each statement.
%
In constraints, we use $\oriset{l}$ to denote the set of origins of a location
$l$ and $\capset{l}$ for the set of capabilities.
%
In addition, $\loc{p}$ denotes the location referred to by a path $p$, defined
as follows, where $\struct{e}$ is the struct type of $e$:
%
\[ \small \loc{x} = x \qquad \loc{e.f} = \struct{e}.f \]
%
Assigning a stream-creating expression to a path generates constraints regarding
the origins based on the assigned value:
%
\[
  %
  \small
  %
  \begin{array}{@{}rcl@{}}
  %
    p={\tt stdin}  &\rightsquigarrow&\ostdin \in\oriset{\loc{p}} \\
  %
    p={\tt stdout} &\rightsquigarrow&\ostdout\in\oriset{\loc{p}} \\
  %
    p={\tt stderr} &\rightsquigarrow&\ostderr\in\oriset{\loc{p}} \\
  %
    p={\tt fopen}()&\rightsquigarrow&\ofile  \in\oriset{\loc{p}} \\
  %
    p={\tt popen}()&\rightsquigarrow&\opipe  \in\oriset{\loc{p}} \\
  %
  \end{array}
  %
\]
%
Passing a path to an API function generates constraints regarding the
capabilities based on the called function:
%
\[
  %
  \small
  %
  \begin{array}{@{}rcl@{}}
  %
    {\tt fread}(p) &\rightsquigarrow&\cread   \in\capset{\loc{p}} \\
  %
    {\tt fscanf}(p)&\rightsquigarrow&\cbufread\in\capset{\loc{p}} \\
  %
    {\tt fwrite}(p)&\rightsquigarrow&\cwrite  \in\capset{\loc{p}} \\
  %
    {\tt fseek}(p) &\rightsquigarrow&\cseek   \in\capset{\loc{p}} \\
  %
    {\tt fclose}(p)&\rightsquigarrow&\cclose  \in\capset{\loc{p}} \\
  %
    {\tt pclose}(p)&\rightsquigarrow&\cclose  \in\capset{\loc{p}} \\
  %
  \end{array}
  %
\]
%
The full list of capabilities required by each API function is provided in the
supplementary material~\cite{supp}.
%
The most interesting case is assignment from one path to another:
%
\[
  %
  \small
  %
  p_1\!=\!p_2\rightsquigarrow
  %
  \oriset{\loc{p_2}}\!\subseteq\!\oriset{\loc{p_1}},
  %
  \capset{\loc{p_1}}\!\subseteq\!\capset{\loc{p_2}}
  %
\]
%
Since all streams stored at $p_2$ are also stored at $p_1$, each member of
$\oriset{\loc{p_2}}$ must belong to $\oriset{\loc{p_1}}$.
%
On the other hand, streams stored at $p_2$ must provide all the capabilities
required by those stored at $p_1$, so $\capset{\loc{p_2}}$ must contain each
member of $\capset{\loc{p_1}}$.
%
This is analogous to subtyping in object-oriented languages:
%
a subtype has more capabilities (methods) than a supertype, and a subtype can be
assigned to a supertype, but not vice versa.

Consider the following example with locations \C{x}, \C{y}, and \C{z}:

\begin{lstlisting}[language=C]
x = fopen(); fseek(x); y = stdin;
if (b) { z = x; } else { z = y; } fread(z);
\end{lstlisting}

\noindent
%
The analysis collects the following constraints for this program:
%
$\ofile\in\oriset{\C{x}}$,
%
$\cseek\in\capset{\C{x}}$,
%
$\ostdin\in\oriset{\C{y}}$,
%
$\oriset{\C{x}}\subseteq\oriset{\C{z}}$,
%
$\capset{\C{z}}\subseteq\capset{\C{x}}$,
%
$\oriset{\C{y}}\subseteq\oriset{\C{z}}$,
%
$\capset{\C{z}}\subseteq\capset{\C{y}}$,
%
and $\cread\in\capset{\C{z}}$.
%
Solving these constraints results in the following sets:
%
$\oriset{\C{x}}=\{\ofile\}$, $\capset{\C{x}}=\{\cseek,\cread\}$,
%
$\oriset{\C{y}}=\{\ostdin\}$, $\capset{\C{y}}=\{\cread\}$,
%
$\oriset{\C{z}}=\{\ofile,\ostdin\}$, and $\capset{\C{z}}=\{\cread\}$.

Since all constraints have the form of set membership or a subset relation, we
can solve them using the well-known cubic algorithm~\cite{spa}.
%
It constructs a graph representing the constraints and then updates it to
propagate the constraints across nodes by following the edges.
%
The algorithm has a time complexity of $O(n^3)$, where $n$ is the number of
locations.
%
In theory, if a target program contains a large number of locations that store
streams, the analysis may not scale well.
%
However, in practice, even large programs use streams in only a portion of the
codebase, and the number of analyzed locations remains relatively small, making
this analysis practically efficient.

\subsection{Error Source Analysis}
\label{sec:analysis:error}

This analysis aims to find all I/O API function calls that serve as the sources
of errors checked by each call to \C{feof} or \C{ferror}.
%
For each error-checking call, we independently perform a dataflow analysis to
identify the sources by starting from the call statement and traversing the
program backward.
%
Before describing the analysis algorithm in detail, we first present examples to
illustrate the intuition behind the analysis.

First, the analysis must handle interprocedural dataflow because error checking
and its sources may not reside in the same function.
%
For example, sources can appear in callees:

\begin{lstlisting}[language=C]
void foo() { bar(f); if (ferror(f)) { .. } }
void bar(FILE *g) { fread(g); }
\end{lstlisting}

\noindent
%
In this example, \C{foo} calls \C{bar}, which performs a read operation on the
stream, and then checks for an error.
%
Thus, the checking occurs in \C{foo}, while the source is the \C{fread} call in
\C{bar}.

It is also possible for sources to be in callers:

\begin{lstlisting}[language=C]
void foo(FILE *f) { if (ferror(f)) { .. } }
void bar() { FILE *g = ..; fread(g); foo(g); }
\end{lstlisting}

\noindent
%
Here, \C{foo} can be viewed as a custom error-checking function, and \C{bar}
calls \C{foo} after reading from the stream.
%
Therefore, the \C{fread} call in \C{bar} is a source for the \C{ferror} call in
\C{foo}.

Second, the analysis can stop moving backward from callees to callers if at
least one source is found.
%
Consider below:

\begin{lstlisting}[language=C]
void foo(FILE *f) { if (cond) { fread(f); }
                    if (ferror(f)) { .. } }
\end{lstlisting}

\noindent
%
This function reaches the \C{ferror} call without performing any I/O operations
if \C{cond} is false.
%
However, the intent of \C{ferror} is presumably to check whether \C{fread}
succeeded, not to detect errors that occurred before entering \C{foo}.
%
When an operation on a stream fails, the next operation on the same stream is
likely to fail again for the same reason.
%
For this reason, C programmers rarely perform operations on a stream whose
indicator may be already set.
%
Therefore, if at least one source is found before reaching the function entry,
our analysis does not continue into the callers.

We now describe the analysis algorithm, shown in \cref{alg:error}.
%
The algorithm \textit{findSources} identifies the sources of an error-checking
call at a label $L_0$.
%
We assume that each statement has a distinct label $L$.
%
For simplicity, we consider only variables and single-parameter functions;
%
the logic can be easily extended to handle struct fields and multiple
parameters.
%
The argument \textit{visitedCallSites} is initially an empty set and is updated
when the algorithm is called recursively.
%
The algorithm consists of two phases:
%
(1) finding sources in the current function and its callees
(lines~\ref{line:error1}--\ref{line:error20}), and
%
(2) finding sources in the callers
(lines~\ref{line:error21}--\ref{line:error31}).

\SetStartEndCondition{ }{}{}
\SetKwProg{Fn}{def}{\string:}{}
\SetKwFor{While}{while}{:}{fintq}
\SetKwFor{For}{for}{\string:}{}
\SetKwIF{If}{ElseIf}{Else}{if}{:}{else if}{else:}{}
\AlgoDontDisplayBlockMarkers\SetAlgoNoEnd\SetAlgoNoLine

\makeatletter
\patchcmd{\@algocf@finish}
  {\@algoskip} % original
  {\@algoskip \vspace{-2em}}
  {}{}
\makeatother

\begin{algorithm}[t]
  \caption{Error source analysis}
  \label{alg:error}
  \setstretch{0.8}
  \footnotesize
  \Fn{findSources($L_0$, visitedCallSites)}{
    $x_0:=$ the argument of the function call at $L_0$\label{line:error1}\;
    \embox{worklist} $=\{(x_0,L_0)\}$; \embox{visited} $=\emptyset$; \embox{sources} $=\emptyset$\label{line:error2}\;
    \While{worklist $\not=\emptyset$\label{line:error3}}{
      $(x,L):=$ \embox{worklist}.\embox{pop}()\label{line:error4}\;
      \lIf{$(x,L)\in$ visited}{
        {\bf continue}\label{line:error6}
      }
      \embox{visited} $:=$ \embox{visited} $\cup\ \{(x,L)\}$\label{line:error7}\;
      $s:=$ the statement at $L$\label{line:error8}\;
      \If{$s=$ \embox{func}$(x)$ for some \embox{func}\label{line:error9}}{
        \If{\embox{func} is a source\label{line:error10}}{
          \embox{sources} $:=$ \embox{sources} $\cup\ \{L\}$\label{line:error11}\;
          {\bf continue}\label{line:error12}\;
        }
        \ElseIf{\embox{func} is a user-defined function\label{line:error13}}{
          $x':=$ the parameter of \embox{func}\label{line:error14}\;
          $L':=$ the label of the return statement of \embox{func}\label{line:error15}\;
          \embox{worklist} $:=$ \embox{worklist} $\cup\ \{(x',L')\}$\label{line:error16}\;
        }
        \lElseIf{\embox{func} is a function ptr}{
          \Return{$\emptyset$}\label{line:error18}
        }
      }
      \For{$L'\leftarrow$ the predecessors of $L$\label{line:error19}}{
        \embox{worklist} $:=$ \embox{worklist} $\cup\ \{(x,L')\}$\label{line:error20}\;
      }
    }
    \embox{func} $:=$ the function containing $L_0$\label{line:error21}\;
    \If{sources $=\emptyset\land x_0$ {\scriptsize is a param} $\land\ $\embox{func}
    {\scriptsize is not used as a function ptr}\label{line:error22}}{
      \embox{visitedCallSites} $:=$ \embox{visitedCallSites} $\cup\ \{L_0\}$\label{line:error23}\;
      \For{$L'\leftarrow$ the callsites of \embox{func}\label{line:error24}}{
        \If{$L'\not\in$ visitedCallSites\label{line:error26}}{
          \embox{sources}$'$ $:=$ \embox{findSources}($L'$, \embox{visitedCallSites})\label{line:error27}\;
          \lIf{sources$'=\emptyset$}{
            \Return{$\emptyset$}\label{line:error29}
          }
          \embox{sources} $:=$ \embox{sources} $\cup$ \embox{sources}$'$\label{line:error30};
        }
      }
    }
    \Return{sources}\label{line:error31}\;
  }
\end{algorithm}

In the first phase, we start by initializing \textit{worklist}, which contains
pairs of a variable and a label to be visited, \textit{visited}, which stores
already-visited pairs, and \textit{sources}, which collects the source labels
(lines~\ref{line:error1}--\ref{line:error2}).
%
Then, we iteratively visit each label until \textit{worklist} becomes empty
(line~\ref{line:error3}).
%
At each iteration, we pop a pair from \textit{worklist}, check whether it has
already been visited, and if not, add it to \textit{visited}
(lines~\ref{line:error4}--\ref{line:error7}).

If the statement at the current label is a function call on the variable
(lines~\ref{line:error8}--\ref{line:error9}), we handle three cases:
%
\begin{itemize}
    %
  \item
    %
    If the callee is an API function that performs a failable operation, we add
    the label to \textit{sources} and do not proceed to the preceding
    statements, as the source along this execution path has now been found
    (lines~\ref{line:error10}--\ref{line:error12}).
    %
  \item
    %
    If the callee is a user-defined function, we add the parameter name and the
    callee's return label to \textit{worklist} to search for sources inside the
    callee (lines~\ref{line:error13}--\ref{line:error16}).
    %
  \item
    %
    If the callee is a function pointer, the algorithm returns an empty set,
    signifying failure to find sources (line~\ref{line:error18}).
    %
    We do not propagate errors across function pointer calls because this
    requires changing the parameter/return types of all functions that can be
    assigned to this function pointer variable, even if they are unrelated to
    I/O.
    %
    Failing to find sources makes the stream unsupported (see
    \cref{sec:analysis:unsupported}).
    %
    %This is to avoid excessive modifications in functions unrelated to I/O.
    %
    %When transforming the program to propagate errors, we must change the
    %parameter and return types of the involved functions to pass error values.
    %%
    %If a source is in a function called via a function pointer, changing the
    %function's type also changes the function pointer's type, which requires
    %changes in all other functions that can be assigned to the same function
    %pointer variable.
    %%
    %Since many of these functions may not perform I/O, modifying all of them
    %would be undesirable.
    %%
    %Thus, we do not handle function pointer calls.
    %
\end{itemize}
%
Finally, we add all the predecessors to \textit{worklist} to continue moving
backward, determining the predecessors of $L$ from the function's control flow
graph (lines~\ref{line:error19}--\ref{line:error20}).
%
We store $x$ in \textit{worklist} because we track both the visited label and
the variable under consideration when visiting a label.
%
This is necessary since multiple variables may store a stream.

The second phase begins by deciding whether we need to proceed to the callers
(lines~\ref{line:error21}--\ref{line:error22}).
%
If no sources have been found yet and $x_0$ is a parameter, we proceed to
analyze the callers.
%
However, if the current function is used as a function pointer, we return the
empty set due to the aforementioned design decision regarding function pointers.

For each call site of the function, we recursively call \textit{findSources}
(lines~\ref{line:error23}--\ref{line:error27}).
%
During this process, we track the set of visited call sites to ensure
termination, even when analyzing recursive functions.
%
If source finding fails at some call site, the overall process is considered to
have failed, and the algorithm returns an empty set
(line~\ref{line:error29}).
%
Otherwise, we accumulate all discovered sources and return them
(lines~\ref{line:error30}--\ref{line:error31}).

We apply \textit{findSources} to each error-checking call to identify its
sources.
%
With memoization of the results from each run of the algorithm, the overall
analysis has a time complexity of $O(n^2)$, where $n$ is the size of the
program, as \textit{findSources} is invoked at $O(n)$ call sites, and each
visits up to $O(n)$ labels.
%
In practice, the analysis often terminates in linear time because sources are
typically located close to the error-checking calls.

\subsection{Unsupported Streams}
\label{sec:analysis:unsupported}

We now discuss the reasons for classifying certain locations as unsupported.
%
Such locations retain the \C{FILE*} type even after the transformation.
%
The majority of the reasons for being unsupported stems from the limited
functionality provided by Rust std compared to libc.
%
We summarize 11 reasons below.

\paragraph{Setbuf}

A stream is passed to the \C{setbuf} or \C{setvbuf} function, which changes the
buffering mode of the stream.
%
However, Rust std does not support changing the buffering mode of standard input
and output to unbuffered.

\paragraph{Ungetc}

A stream is passed to the \C{ungetc} function, which pushes a byte into a
readable stream.
%
However, \C{Read} and \C{BufRead} in Rust std do not provide such functionality.

\paragraph{Freopen}

A stream is passed to the \C{freopen} function, which changes the underlying
object to a specified file.
%
Rust std does not support this behavior.

\paragraph{Improper Capabilities}

The capabilities required of a stream do not match its possible origins.
%
For example, in libc, seeking on standard input may succeed at run time if the
input has been redirected to a file;
%
in Rust std, however, \C{Stdin} never implements the \C{Seek} trait, regardless
of redirection.

\paragraph{Cyclic Close Capability}

The capabilities of multiple streams form a \emph{cyclic} dependency while
including $\cclose$. The following code is an example:

\vspace{-0.25em}
\begin{lstlisting}[language=C]
x = y; .. y = x; .. fclose(x);
\end{lstlisting}
\vspace{-0.25em}

\noindent
%
Here, the capabilities of \C{x} and \C{y} form a cycle, as they are assigned to
each other, and $\cclose$ is required due to the \C{fclose} call.
%
While this works in C, Rust's ownership type system allows each stream to be
owned by a single variable, preventing this form of cyclic ownership.

\paragraph{Comparison}

Two \C{FILE*} pointers are compared using \CT{==} or \CT{!=}.
%
Rust std does not support comparing streams.

\paragraph{Casts}

A stream is cast to or from a \C{void} pointer or an integer.
%
Since Rust std's streams are not pointers, they cannot be cast to or from
\C{void} pointers or integers.

\paragraph{Variadic}

A stream is passed as a variadic argument.
%
Variadic arguments are not type-checked, so it is meaningless to assign Rust
std's types to them.

\paragraph{API Functions as Pointers}

A stream is passed to a function pointer call that may refer to an I/O API
function from libc.
%
API functions such as \C{feof} and \C{ferror} cannot be represented as function
pointers using Rust std.

\paragraph{Error Sources Unfound}

A stream is passed to \C{feof} or \C{ferror}, but error source analysis failed
to find the sources.

\paragraph{Non-POSIX}

The stream is passed to a non-POSIX API function, such as \C{\_\_freading} in
GNU libc.
%
Non-POSIX APIs are outside the scope of this work.

\vspace{0.5em}

The unsupported-ness of a location propagates through assignments,
\emph{bidirectionally}, from the right value to the left value and vice versa.
%
For example, in the following code, if \C{y} is unsupported, then not only \C{x}
but also \C{z} become unsupported:

\begin{lstlisting}[language=C]
if (..) { x = y; } else { x = z; }
\end{lstlisting}

\noindent
%
This is because we cannot convert a libc stream to or from a Rust std stream at
run time.
%
If \C{z} is a Rust std stream while \C{y} is a libc stream, they cannot be
assigned to the same variable.

Since the unsupported-ness of the left value and right value of an assignment is
the same, we can compute the complete set of unsupported locations using the
unification algorithm based on the union-find data
structure~\cite{galler1964improved}.
%
This runs in $O(n \cdot \alpha(n))$ time complexity, where $\alpha$ is the
inverse Ackermann function.
